% lang/es.tex
% -----------------------------------------------------------------------
% Todas las cadenas de texto del cuaderno, en un solo sitio, como en
% "Aprendo a leer": cambiar una palabra aquí la cambia en las 260 páginas
% a la vez, sin tocar el contenido generado.
% -----------------------------------------------------------------------

% Total de días del cuaderno. tools/gen_numeros.py --check comprueba que
% coincide con su propio TOTAL_DIAS.
\newcommand{\totaldias}{260}

% La cabecera de cada día (\cabeceraDia, preamble.tex).
\newcommand{\lblFecha}{Fecha}
\newcommand{\lblNotas}{Notas}
\newcommand{\lblDia}{Día}
\newcommand{\lblSemana}{Semana}
\newcommand{\lblTrimestre}{Trimestre}
\newcommand{\lblHecho}{Hecho}
\newcommand{\lblHechoSola}{sola}
\newcommand{\lblHechoConAyuda}{con ayuda}
\newcommand{\lblHechoConDificultad}{con dificultad}

% La caja de arriba, la misma todos los días.
\newcommand{\lblNumeroDeHoy}{El número de hoy}

% Los títulos de las cajas de actividad.
\newcommand{\lblTraza}{Traza}
\newcommand{\lblColorea}{Colorea}
\newcommand{\lblRodea}{Rodea}
\newcommand{\lblCuenta}{Cuenta}
\newcommand{\lblBusca}{Busca}
\newcommand{\lblUne}{Une}
\newcommand{\lblCompleta}{Completa}
\newcommand{\lblDecena}{Diez y más}
\newcommand{\lblCompara}{Compara}
\newcommand{\lblVecinos}{Antes y después}
\newcommand{\lblRecta}{La recta}
\newcommand{\lblOrdena}{Ordena}
\newcommand{\lblJunta}{Junta}
\newcommand{\lblSuma}{Suma}
\newcommand{\lblResta}{Resta}
\newcommand{\lblParte}{Parte en dos}
\newcommand{\lblDiez}{Hasta el 10}
\newcommand{\lblProblema}{Problema}
\newcommand{\lblBloques}{Decenas y unidades}
% "Number names" solo sale en "First Numbers" (lang/en.tex), pero su caja
% es la misma en los dos cuadernos (preamble.tex).
\newcommand{\lblNombres}{Los nombres de los números}
\newcommand{\lblTabla}{La tabla del 100}
\newcommand{\lblDinero}{El dinero}
\newcommand{\lblHora}{La hora}
\newcommand{\lblMide}{Mide}
\newcommand{\lblDibuja}{Dibuja}
\newcommand{\lblRepasa}{Repasa}

% Las instrucciones: las lee un adulto en voz alta (van pequeñas y en
% gris, como en "Aprendo a leer"). Las de "colorea", "rodea", "cuenta",
% "une" y "completa" las compone tools/gen_numeros.py con el número y las
% cosas de cada día ("Colorea 2 pelotas.").
\newcommand{\lblInstruccionTraza}{Repasa los puntos con el dedo, y luego con un lápiz. Después, escríbelo tú en la línea.}
\newcommand{\lblInstruccionBusca}{Rodea cada uno que encuentres. ¿Cuántos hay?}

% La tabla de las decenas (D) y las unidades (U) de la caja de arriba.
\newcommand{\lblDecenasCorto}{D}
\newcommand{\lblUnidadesCorto}{U}

% La portada.
\newcommand{\lblTituloPortada}{Aprendo los números}
\newcommand{\lblSubtituloPortada}{Un número nuevo cada semana}
\newcommand{\lblSubtituloPortadaMetodo}{Cuaderno diario para contar, trazar los números y empezar a sumar -- con la familia de \emph{Aprendo a leer}}
\newcommand{\lblPortadaNombre}{Nombre}
\newcommand{\lblPortadaCurso}{Curso}

% La medalla del final de cada trimestre (T1-T3) y el cartel de cada
% diez páginas.
\newcommand{\lblMedalla}[1]{¡Medalla de #1!}
\newcommand{\lblMedallaAnimo}{¡Sigue así, \rule{55mm}{0.4pt}!}
\newcommand{\lblPaginasHechas}[1]{¡#1 páginas hechas!}

% La clave de respuestas.
\newcommand{\lblClave}{Clave de respuestas}
